\section{Local construction of unramified and \'etale morphisms}
\label{I.7}
% original p. 10

\begin{proposition}
\label{I.7.1}
Let $A$ be a noetherian ring, $B$ a finite algebra over~$A$,
$u$ a generator of~$B$ over~$A$, $F\in A[t]$ such that $F(u)=0$
($F$ is not assumed monic), $u^\prime=F^\prime(u)$ (where
$F^\prime$ is the derivative polynomial), $\goth{q}$ a prime
ideal of~$B$ not containing~$u^\prime$, and $\goth{p}$ its
contraction to~$A$. Then $B_\goth{q}$ is net over~$A_\goth{p}$.
\end{proposition}

In other words, setting $Y=\Spec(A)$, $X=\Spec(B)$,
$X_{u^\prime}=\Spec(B_{u^\prime})$, $X_{u^\prime}$ is unramified
over~$Y$. The statement follows from the following, more precise one:

\begin{corollary}
\label{I.7.2}
The different ideal of~$B/A$ contains $u^\prime B$, and is equal
to it if the natural homomorphism $A[t]/FA[t]\to B$ (sending
$t$ to~$u$) is an isomorphism.
\end{corollary}

Let $J$ be the kernel of the homomorphism $C=A[t]\to B$; this
kernel contains $FA[t]$, and is equal to it in the second case
considered in~\ref{I.7.2}. Since the homomorphism is surjective,
$\Omega^1_{B/A}$ is identified with the quotient of~$\Omega^1_{C/A}$
by the submodule generated by $J\Omega^1_{C/A}$ and $d(J)$
(the definition of the homomorphism~$d$, and the calculation
of~$\Omega^1$ for a polynomial algebra, should have been made
explicit in \No\ref{I.1}). Identifying $\Omega^1_{C/A}$ with~$C$
by means of the basis~$dt$, one finds $B/B\cdot J^\prime$;
thus the different is generated by the set~$J^\prime$ of images
in~$B$ of the derivatives of the $G\in J$ (and it suffices to
take $G$ generating~$J$). Since $F\in J$, \resp $F$ is a generator
of~$J$, we are done. (N.B.\ One should make~\ref{I.7.2} a
proposition and~\ref{I.7.1} a corollary.) One obtains:

\begin{corollary}
\label{I.7.3}
Under the conditions of~\ref{I.7.1}, assuming that $F$ is monic
and that $A[t]/FA[t]\to B$ is an isomorphism, for $B_\goth{q}$
to be \'etale over~$A_\goth{p}$, it is necessary and sufficient
that~$\goth{q}$ not contain~$u^\prime$.
\end{corollary}

Indeed, since $B$ is flat over~$A$, \'etale is equivalent to
net, and one can apply~\ref{I.7.2}.

\begin{corollary}
\label{I.7.4}
Under the conditions of~\ref{I.7.3}, for $B$ to be \'etale
over~$A$, it is necessary and sufficient that~$u^\prime$ be
invertible, or equivalently that the ideal generated by $F$,
$F^\prime$ in $A[t]$ be the unit ideal.
\end{corollary}

The last criterion follows from the first and Nakayama (in~$B$).

A monic polynomial $F\in A[t]$ having the property stated in
Corollary~\ref{I.7.4}
% original p. 11
is called a \emph{separable polynomial} (if $F$ is not monic,
one should at least require its leading coefficient to be
invertible; when $A$ is a field, one recovers the usual definition).

\begin{corollary}
\label{I.7.5}
Let $B$ be a \emph{finite} algebra over the \emph{local} ring~$A$.
Suppose that $K(A)$ is infinite or that $B$ is local. Let~$n$
be the rank of~$L=B\otimes_A K(A)$ over~$K(A)=k$. For $B$ to
be net (\resp \'etale) over~$A$, it is necessary and sufficient
that $B$ be isomorphic to a quotient of (\resp isomorphic to)
$A[t]/FA[t]$, where $F$ is a monic separable polynomial which
one may assume to be (\resp which is necessarily) of degree~$n$.
\end{corollary}

Only necessity needs to be proved. Suppose that $B$ is net
over~$A$, so that $L$ is separable over~$k$. It then follows
from the hypothesis that $L/k$ has a generator~$\xi$, so the
$\xi^i$ ($0\leq i<n$) form a basis of~$L$ over~$k$. Let
$u\in B$ lift~$\xi$; then, by Nakayama, the $u^i$ ($0\leq i<n$)
generate the $A$-module~$B$ (\resp form a basis of it).
In particular, one can find a monic polynomial $F\in A[t]$
such that $F(u)=0$, and $B$ will be isomorphic to a quotient
of (\resp isomorphic to) $A[t]/FA[t]$. Finally, by~\ref{I.7.4}
applied to $L/k$, $F$ and~$F^\prime$ generate $A[t]$ modulo
$\goth{m}A[t]$; hence (by Nakayama in $A[t]/FA[t]$) $F$ and
$F^\prime$ generate~$A[t]$, and we are done.

\begin{theorem}
\label{I.7.6}
Let $A$ be a local ring, and $A\to\cal{O}$ a local homomorphism
such that~$\cal{O}$ is isomorphic to a localization of an algebra
of finite type over~$A$. Suppose that $\cal{O}$ is \emph{net}
over~$A$. Then one can find an $A$-algebra $B$, integral over~$A$,
a maximal ideal $\goth{n}$ of~$B$, a generator $u$ of~$B$
over~$A$, and a monic polynomial $F\in A[t]$, such that
$\goth{n}\not\ni F^\prime(u)$ and $\cal{O}$ is isomorphic
(as an $A$-algebra) to~$B_{\goth{n}}$. If $\cal{O}$ is \'etale
over~$A$, one can take $B=A[t]/FA[t]$.
\end{theorem}

(Of course, these conditions are also sufficient\ldots)

Let us first point out the pleasant corollaries:

\begin{corollary}
\label{I.7.7}
For $\cal{O}$ to be net over~$A$, it is necessary and sufficient
that $\cal{O}$ be isomorphic to the quotient of an analogous
algebra which is \emph{\'etale} over~$A$.
\end{corollary}

Indeed, one takes $\cal{O}^\prime=B^\prime_{\goth{n}^\prime}$,
where $B^\prime=A[t]/FA[t]$ and $\goth{n}^\prime$ is the inverse
image of~$\goth{n}$ in~$B^\prime$.

\begin{corollary}
\label{I.7.8}
Let
% original p. 12
$f\colon X\to Y$ be a morphism of finite type, $x\in X$.
For $f$ to be net at~$x$, it is necessary and sufficient that
there exist an open neighborhood $U$ of~$x$ such that $f|U$
factors as $U\to X^\prime\to Y$, where the first arrow is
a closed immersion and the second an \'etale morphism.
\end{corollary}

This is a mere translation of~\ref{I.7.7}.

Let us show how the jargon of~\ref{I.7.6} follows from the
main statement: indeed, by~\ref{I.7.7} there is an epimorphism
$\cal{O}^\prime\to\cal{O}$, where $\cal{O}$ has the desired
properties; but since $\cal{O}^\prime$ and~$\cal{O}$ are \'etale
over~$A$, the morphism $\cal{O}^\prime\to\cal{O}$ is \'etale
by~\ref{I.4.8}, hence an isomorphism.

\subsubsection*{Proof of~\ref{I.7.6}}
It follows a proof from the Chevalley seminar. By the
\emph{Main Theorem}, one has $\cal{O}=B_{\goth{n}}$, where
$B$ is a finite algebra over~$A$ and $\goth{n}$ is a maximal
ideal of it. Then $B/\goth{n}=K(\cal{O})$ is a separable,
hence monogenic, extension of~$k$; if $\goth{n}_i$
($1\leq i\leq r$) are the maximal ideals of~$B$ distinct
from~$\goth{n}$, there is therefore an element $u$ of~$B$
which belongs to all the~$\goth{n}_i$, and whose image in
$B/\goth{n}$ is a generator of it. Now
$B/\goth{n}=B_\goth{n}/\goth{n}B_\goth{n}
=B_\goth{n}/\goth{m}B_\goth{n}$ (where $\goth{m}$ is the
maximal ideal of~$A$). Let us admit for a moment the following

\begin{lemma}
\label{I.7.9}
Let $A$ be a local ring, $B$ a finite algebra over~$A$,
$\goth{n}$ a maximal ideal of~$B$, and $u$ an element of~$B$
whose image in $B_\goth{n}/\goth{m}B_\goth{n}$ generates it
as an algebra over~$k=A/\goth{m}$, and which belongs to every
maximal ideal of~$B$ distinct from~$\goth{n}$. Let
$B^\prime=B[u]$, $\goth{n}^\prime=\goth{n}B^\prime$.
Then the canonical homomorphism
$B^\prime_{\goth{n}^\prime}\to B_\goth{n}$ is an isomorphism.
\end{lemma}

\begin{lemma}
\label{I.7.10}
(Should have appeared as a corollary to~\ref{I.7.1},
before~\ref{I.7.5}, which it implies.) Let $B$ be a finite
algebra over~$A$ generated by an element~$u$, and let
$\goth{n}$ be a maximal ideal of~$B$ such that $B_\goth{n}$
is unramified over~$A$. Then there exists a monic polynomial
$F\in A[t]$ such that $F(u)=0$ and $F^\prime(u)\notin\goth{n}$.
\end{lemma}

Indeed, let $n$ be the rank of the $k$-algebra $L=B\otimes_A k$;
by Nakayama there exists a monic polynomial of degree~$n$
in $A[t]$ such that $F(u)=0$. Let $f$ be the polynomial obtained
from~$F$ by reduction mod~$\goth{m}$; then $L$ is $k$-isomorphic
to $k[t]/fk[t]$, so by~\ref{I.7.3}, $f^\prime(\xi)$ does not
belong to the maximal ideal of~$L$ corresponding to~$\goth{n}$
($\xi$ denoting the image of~$t$ in~$L$, \ie the image of~$u$
in~$L$). Since $f^\prime(\xi)$ is the image of~$F^\prime(u)$,
we are done.

The
% original p. 13
theorem~\ref{I.7.6} now follows by combining~\ref{I.7.9}
and~\ref{I.7.10}. It remains to prove~\ref{I.7.9}. Set
$S^\prime=B^\prime-\goth{n}^\prime$, so
$B^\prime{S^\prime}^{-1}=B^\prime_{\goth{n}^\prime}$.
\[
\begin{array}{ccccc}
B&\longrightarrow&B{S^\prime}^{-1}&\longrightarrow&BS^{-1} =
B_\goth{n}\\ \Big\uparrow&&\Big\uparrow&&\\
B^\prime&\longrightarrow&B^\prime{S^\prime}^{-1}
\rlap{\hbox{$\displaystyle \ = B^\prime_{\goth{n}^\prime}$}}&&
\end{array}
\]
Similarly, let $S=B-\goth{n}$, so $BS^{-1}=B_\goth{n}$.
There is therefore a natural homomorphism
$B{S^\prime}^{-1}\to BS^{-1}=B_\goth{n}$; let us prove that
it is an isomorphism, \ie that the elements of~$S$ are invertible
in $B{S^\prime}^{-1}$, \ie that every maximal ideal $\goth{p}$
of the latter is disjoint from~$S$, \ie contracts to~$\goth{n}$
in~$B$. Indeed, since $B{S^\prime}^{-1}$ is finite over
$B^\prime{S^\prime}^{-1}=B^\prime_{\goth{n}^\prime}$,
$\goth{p}$ contracts to the unique maximal ideal
$\goth{n}^\prime B_{\goth{n}^\prime}$ of
$B^\prime_{\goth{n}^\prime}$, hence contracts to the maximal
ideal $\goth{n}^\prime$ of~$B^\prime$; since $B$ is finite
over~$B^\prime$, the ideal $\goth{q}$ of~$B$ induced
by~$\goth{p}$, lying over~$\goth{n}^\prime$, is necessarily
maximal, and does not contain~$u$, hence is identical
to~$\goth{n}$. (We have just used the fact that $u$ belongs
to every maximal ideal of~$B$ distinct from~$\goth{n}$.)
Let us now prove that $B{S^\prime}^{-1}$ equals
$B^\prime{S^\prime}^{-1}$: since it is finite over the latter,
Nakayama reduces one to proving equality mod
$\goth{n}^\prime B{S^\prime}^{-1}$, and a fortiori it suffices
to prove equality mod $\goth{m}B{S^\prime}^{-1}$; now
$B{S^\prime}^{-1}/\goth{m}B{S^\prime}^{-1}
=B_\goth{n}/\goth{m}B_\goth{n}$ is generated over~$k$ by~$u$
(here one uses the other property of~$u$), so the image
of~$B^\prime$ (and a fortiori of $B^\prime{S^\prime}^{-1}$)
in it is the whole ring (being a subring containing~$k$
and the image of~$u$).

\begin{remark}
It should be possible to state Theorem~\ref{I.7.6} for a ring
$\cal{O}$ which is merely semilocal, so as to include~\ref{I.7.5}
as well: one would assume that $\cal{O}/\goth{m}\cal{O}$ is a
\emph{monogenic} $k$-algebra; one could therefore find a
$u\in B$ whose image in $B/\goth{m}B$ is a generator, and
which belongs to all the maximal ideals of~$B$ not coming
from~$\cal{O}$. Lemmas~\ref{I.7.9} and~\ref{I.7.10} should
adapt without difficulty. More generally, \ldots
\end{remark}
